% Generated by https://pasdoc.github.io/PasDoc 0.16.0
\documentclass{report}
\usepackage{hyperref}
% WARNING: THIS SHOULD BE MODIFIED DEPENDING ON THE LETTER/A4 SIZE
\oddsidemargin 0cm
\evensidemargin 0cm
\marginparsep 0cm
\marginparwidth 0cm
\parindent 0cm
\setlength{\textwidth}{\paperwidth}
\addtolength{\textwidth}{-2in}


% Conditional define to determine if pdf output is used
\newif\ifpdf
\ifx\pdfoutput\undefined
\pdffalse
\else
\pdfoutput=1
\pdftrue
\fi

\ifpdf
  \usepackage[pdftex]{graphicx}
\else
  \usepackage[dvips]{graphicx}
\fi

% Write Document information for pdflatex/pdftex
\ifpdf
\pdfinfo{
 /Author     (Pasdoc)
 /Title      ()
}
\fi


% definitons for warning and note tag
\usepackage[most]{tcolorbox}
\newtcolorbox{tcbwarning}{
 breakable,
 enhanced jigsaw,
 top=0pt,
 bottom=0pt,
 titlerule=0pt,
 bottomtitle=0pt,
 rightrule=0pt,
 toprule=0pt,
 bottomrule=0pt,
 colback=white,
 arc=0pt,
 outer arc=0pt,
 title style={white},
 fonttitle=\color{black}\bfseries,
 left=8pt,
 colframe=red,
 title={Warning:},
}
\newtcolorbox{tcbnote}{
 breakable,
 enhanced jigsaw,
 top=0pt,
 bottom=0pt,
 titlerule=0pt,
 bottomtitle=0pt,
 rightrule=0pt,
 toprule=0pt,
 bottomrule=0pt,
 colback=white,
 arc=0pt,
 outer arc=0pt,
 title style={white},
 fonttitle=\color{black}\bfseries,
 left=8pt,
 colframe=yellow,
 title={Note:},
}

\begin{document}
\label{toc}\tableofcontents
\newpage
% special variable used for calculating some widths.
\newlength{\tmplength}
\chapter{Unit legacey.random}
\label{legacey.random}
\index{legacey.random}
\section{Description}
A very simple random number generator for Blaise pascal. Not Needed for Delphi or FPC compilers
\section{Overview}
\begin{description}
\item[\texttt{random}]
\end{description}
\section{Functions and Procedures}
\ifpdf
\subsection*{\large{\textbf{random}}\normalsize\hspace{1ex}\hrulefill}
\else
\subsection*{random}
\fi
\label{legacey.random-random}
\index{random}
\begin{list}{}{
\settowidth{\tmplength}{\textbf{Description}}
\setlength{\itemindent}{0cm}
\setlength{\listparindent}{0cm}
\setlength{\leftmargin}{\evensidemargin}
\addtolength{\leftmargin}{\tmplength}
\settowidth{\labelsep}{X}
\addtolength{\leftmargin}{\labelsep}
\setlength{\labelwidth}{\tmplength}
}
\item[\textbf{Declaration}\hfill]
\ifpdf
\begin{flushleft}
\fi
\begin{ttfamily}
function random: Double;\end{ttfamily}

\ifpdf
\end{flushleft}
\fi

\par
\item[\textbf{Description}]
Returns a random double from 0 to 1.0

\end{list}
\section{Variables}
\ifpdf
\subsection*{\large{\textbf{RandSeed64}}\normalsize\hspace{1ex}\hrulefill}
\else
\subsection*{RandSeed64}
\fi
\label{legacey.random-RandSeed64}
\index{RandSeed64}
\begin{list}{}{
\settowidth{\tmplength}{\textbf{Description}}
\setlength{\itemindent}{0cm}
\setlength{\listparindent}{0cm}
\setlength{\leftmargin}{\evensidemargin}
\addtolength{\leftmargin}{\tmplength}
\settowidth{\labelsep}{X}
\addtolength{\leftmargin}{\labelsep}
\setlength{\labelwidth}{\tmplength}
}
\item[\textbf{Declaration}\hfill]
\ifpdf
\begin{flushleft}
\fi
\begin{ttfamily}
RandSeed64: UInt64 = 88172645463325252;\end{ttfamily}

\ifpdf
\end{flushleft}
\fi

\par
\item[\textbf{Description}]
Random number sequence will always be the same with same seed! If it is a problem, a seed can be used from system clock.

\end{list}
\chapter{Unit morse.base}
\label{morse.base}
\index{morse.base}
\section{Description}
The base unit for the Morse{-}lib. These are what are needed to configure the morse code generator. A default settings is provided called \begin{ttfamily}DefaultMorseSettings\end{ttfamily}(\ref{morse.base-DefaultMorseSettings}). The settings are as follows : \texttt{\\\nopagebreak[3]
WPM~:~18\\\nopagebreak[3]
Farnsworth~:~Disabled\\\nopagebreak[3]
Noise~:~Disabled\\\nopagebreak[3]
Tone~Hz~~:~800~Hz\\\nopagebreak[3]
Volume~:~(TBA)\\
}
\section{Overview}
\begin{description}
\item[\texttt{\begin{ttfamily}TmorseSettings\end{ttfamily} Record}]
\item[\texttt{\begin{ttfamily}TElementTimings\end{ttfamily} Record}]
\end{description}
\section{Classes, Interfaces, Objects and Records}
\ifpdf
\subsection*{\large{\textbf{TmorseSettings Record}}\normalsize\hspace{1ex}\hrulefill}
\else
\subsection*{TmorseSettings Record}
\fi
\label{morse.base.TmorseSettings}
\index{TmorseSettings}
\subsubsection*{\large{\textbf{Description}}\normalsize\hspace{1ex}\hfill}
The settting used for creation of the morse code sounds\subsubsection*{\large{\textbf{Fields}}\normalsize\hspace{1ex}\hfill}
\begin{list}{}{
\settowidth{\tmplength}{\textbf{FarsworthEnabled}}
\setlength{\itemindent}{0cm}
\setlength{\listparindent}{0cm}
\setlength{\leftmargin}{\evensidemargin}
\addtolength{\leftmargin}{\tmplength}
\settowidth{\labelsep}{X}
\addtolength{\leftmargin}{\labelsep}
\setlength{\labelwidth}{\tmplength}
}
\label{morse.base.TmorseSettings-wpm}
\index{wpm}
\item[\textbf{wpm}\hfill]
\ifpdf
\begin{flushleft}
\fi
\begin{ttfamily}
public wpm: Integer;\end{ttfamily}

\ifpdf
\end{flushleft}
\fi


\par Set the WPM for the morse code output\label{morse.base.TmorseSettings-ToneHertz}
\index{ToneHertz}
\item[\textbf{ToneHertz}\hfill]
\ifpdf
\begin{flushleft}
\fi
\begin{ttfamily}
public ToneHertz: Integer;\end{ttfamily}

\ifpdf
\end{flushleft}
\fi


\par Set the tone in hertz for the audio output\label{morse.base.TmorseSettings-morseLevel}
\index{morseLevel}
\item[\textbf{morseLevel}\hfill]
\ifpdf
\begin{flushleft}
\fi
\begin{ttfamily}
public morseLevel: Integer;\end{ttfamily}

\ifpdf
\end{flushleft}
\fi


\par Set the level of the output. A good start would be 70{\%} of max Integer value\label{morse.base.TmorseSettings-FarsworthEnabled}
\index{FarsworthEnabled}
\item[\textbf{FarsworthEnabled}\hfill]
\ifpdf
\begin{flushleft}
\fi
\begin{ttfamily}
public FarsworthEnabled: Boolean;\end{ttfamily}

\ifpdf
\end{flushleft}
\fi


\par Enable Farnsworth Spacing\label{morse.base.TmorseSettings-FarnsworthWPM}
\index{FarnsworthWPM}
\item[\textbf{FarnsworthWPM}\hfill]
\ifpdf
\begin{flushleft}
\fi
\begin{ttfamily}
public FarnsworthWPM: Integer;\end{ttfamily}

\ifpdf
\end{flushleft}
\fi


\par Set Farnsworth Spacing WPM, if enabled, otherwise, no affect\label{morse.base.TmorseSettings-NoiseType}
\index{NoiseType}
\item[\textbf{NoiseType}\hfill]
\ifpdf
\begin{flushleft}
\fi
\begin{ttfamily}
public NoiseType: TnoiseType;\end{ttfamily}

\ifpdf
\end{flushleft}
\fi


\par Set the noise type to off, white, or pink\label{morse.base.TmorseSettings-NoiseLevel}
\index{NoiseLevel}
\item[\textbf{NoiseLevel}\hfill]
\ifpdf
\begin{flushleft}
\fi
\begin{ttfamily}
public NoiseLevel: Integer;\end{ttfamily}

\ifpdf
\end{flushleft}
\fi


\par Set the noise level. It should be the inverse of the morseLevel, 30 {\%} of max Integer value, for example\end{list}
\ifpdf
\subsection*{\large{\textbf{TElementTimings Record}}\normalsize\hspace{1ex}\hrulefill}
\else
\subsection*{TElementTimings Record}
\fi
\label{morse.base.TElementTimings}
\index{TElementTimings}
\subsubsection*{\large{\textbf{Description}}\normalsize\hspace{1ex}\hfill}
Holds all the timing values for the elements\subsubsection*{\large{\textbf{Fields}}\normalsize\hspace{1ex}\hfill}
\begin{list}{}{
\settowidth{\tmplength}{\textbf{interCharacter}}
\setlength{\itemindent}{0cm}
\setlength{\listparindent}{0cm}
\setlength{\leftmargin}{\evensidemargin}
\addtolength{\leftmargin}{\tmplength}
\settowidth{\labelsep}{X}
\addtolength{\leftmargin}{\labelsep}
\setlength{\labelwidth}{\tmplength}
}
\label{morse.base.TElementTimings-dit}
\index{dit}
\item[\textbf{dit}\hfill]
\ifpdf
\begin{flushleft}
\fi
\begin{ttfamily}
public dit: Single;\end{ttfamily}

\ifpdf
\end{flushleft}
\fi


\par The length of a dit in fractional seconds\label{morse.base.TElementTimings-dah}
\index{dah}
\item[\textbf{dah}\hfill]
\ifpdf
\begin{flushleft}
\fi
\begin{ttfamily}
public dah: Single;\end{ttfamily}

\ifpdf
\end{flushleft}
\fi


\par The length of a dah in fractional seconds is dit * 3\label{morse.base.TElementTimings-intraElement}
\index{intraElement}
\item[\textbf{intraElement}\hfill]
\ifpdf
\begin{flushleft}
\fi
\begin{ttfamily}
public intraElement: Single;\end{ttfamily}

\ifpdf
\end{flushleft}
\fi


\par The length of time between elements\label{morse.base.TElementTimings-interCharacter}
\index{interCharacter}
\item[\textbf{interCharacter}\hfill]
\ifpdf
\begin{flushleft}
\fi
\begin{ttfamily}
public interCharacter: Single;\end{ttfamily}

\ifpdf
\end{flushleft}
\fi


\par The length of time between characters. It can change based on Farnsworth settings. Otherwise it is dit * 3\label{morse.base.TElementTimings-interWord}
\index{interWord}
\item[\textbf{interWord}\hfill]
\ifpdf
\begin{flushleft}
\fi
\begin{ttfamily}
public interWord: Single;\end{ttfamily}

\ifpdf
\end{flushleft}
\fi


\par The length of time between words. It can change based on Farnsworth settings. Otherwise it is dit * 6\end{list}
\section{Types}
\ifpdf
\subsection*{\large{\textbf{Tpcm}}\normalsize\hspace{1ex}\hrulefill}
\else
\subsection*{Tpcm}
\fi
\label{morse.base-Tpcm}
\index{Tpcm}
\begin{list}{}{
\settowidth{\tmplength}{\textbf{Description}}
\setlength{\itemindent}{0cm}
\setlength{\listparindent}{0cm}
\setlength{\leftmargin}{\evensidemargin}
\addtolength{\leftmargin}{\tmplength}
\settowidth{\labelsep}{X}
\addtolength{\leftmargin}{\labelsep}
\setlength{\labelwidth}{\tmplength}
}
\item[\textbf{Declaration}\hfill]
\ifpdf
\begin{flushleft}
\fi
\begin{ttfamily}
Tpcm = array of Integer;\end{ttfamily}

\ifpdf
\end{flushleft}
\fi

\par
\item[\textbf{Description}]
Array used for all PCM Data in morse{-}lib

\end{list}
\ifpdf
\subsection*{\large{\textbf{EmorseLibError}}\normalsize\hspace{1ex}\hrulefill}
\else
\subsection*{EmorseLibError}
\fi
\label{morse.base-EmorseLibError}
\index{EmorseLibError}
\begin{list}{}{
\settowidth{\tmplength}{\textbf{Description}}
\setlength{\itemindent}{0cm}
\setlength{\listparindent}{0cm}
\setlength{\leftmargin}{\evensidemargin}
\addtolength{\leftmargin}{\tmplength}
\settowidth{\labelsep}{X}
\addtolength{\leftmargin}{\labelsep}
\setlength{\labelwidth}{\tmplength}
}
\item[\textbf{Declaration}\hfill]
\ifpdf
\begin{flushleft}
\fi
\begin{ttfamily}
EmorseLibError = Exception;\end{ttfamily}

\ifpdf
\end{flushleft}
\fi

\par
\item[\textbf{Description}]
Used for all exceptions in morse{-}lib

\end{list}
\ifpdf
\subsection*{\large{\textbf{TsampleRate}}\normalsize\hspace{1ex}\hrulefill}
\else
\subsection*{TsampleRate}
\fi
\label{morse.base-TsampleRate}
\index{TsampleRate}
\begin{list}{}{
\settowidth{\tmplength}{\textbf{Description}}
\setlength{\itemindent}{0cm}
\setlength{\listparindent}{0cm}
\setlength{\leftmargin}{\evensidemargin}
\addtolength{\leftmargin}{\tmplength}
\settowidth{\labelsep}{X}
\addtolength{\leftmargin}{\labelsep}
\setlength{\labelwidth}{\tmplength}
}
\item[\textbf{Declaration}\hfill]
\ifpdf
\begin{flushleft}
\fi
\begin{ttfamily}
TsampleRate = (...);\end{ttfamily}

\ifpdf
\end{flushleft}
\fi

\par
\item[\textbf{Description}]
sample rates for PCM data\item[\textbf{Values}]
\begin{description}
\item[\texttt{srLow = 8000}] \label{morse.base-srLow}
\index{}
 
\item[\texttt{srMedium = 22050}] \label{morse.base-srMedium}
\index{}
 
\item[\texttt{srCD = 44100}] \label{morse.base-srCD}
\index{}
 
\item[\texttt{srDVD = 480000}] \label{morse.base-srDVD}
\index{}
 
\end{description}


\end{list}
\ifpdf
\subsection*{\large{\textbf{TnoiseType}}\normalsize\hspace{1ex}\hrulefill}
\else
\subsection*{TnoiseType}
\fi
\label{morse.base-TnoiseType}
\index{TnoiseType}
\begin{list}{}{
\settowidth{\tmplength}{\textbf{Description}}
\setlength{\itemindent}{0cm}
\setlength{\listparindent}{0cm}
\setlength{\leftmargin}{\evensidemargin}
\addtolength{\leftmargin}{\tmplength}
\settowidth{\labelsep}{X}
\addtolength{\leftmargin}{\labelsep}
\setlength{\labelwidth}{\tmplength}
}
\item[\textbf{Declaration}\hfill]
\ifpdf
\begin{flushleft}
\fi
\begin{ttfamily}
TnoiseType = (...);\end{ttfamily}

\ifpdf
\end{flushleft}
\fi

\par
\item[\textbf{Description}]
Noise types for audio PCM data (off, white, or pink)\item[\textbf{Values}]
\begin{description}
\item[\texttt{ntOff}] \label{morse.base-ntOff}
\index{}
 
\item[\texttt{ntWhite}] \label{morse.base-ntWhite}
\index{}
 
\item[\texttt{ntPink}] \label{morse.base-ntPink}
\index{}
 
\end{description}


\end{list}
\section{Variables}
\ifpdf
\subsection*{\large{\textbf{DefaultMorseSettings}}\normalsize\hspace{1ex}\hrulefill}
\else
\subsection*{DefaultMorseSettings}
\fi
\label{morse.base-DefaultMorseSettings}
\index{DefaultMorseSettings}
\begin{list}{}{
\settowidth{\tmplength}{\textbf{Description}}
\setlength{\itemindent}{0cm}
\setlength{\listparindent}{0cm}
\setlength{\leftmargin}{\evensidemargin}
\addtolength{\leftmargin}{\tmplength}
\settowidth{\labelsep}{X}
\addtolength{\leftmargin}{\labelsep}
\setlength{\labelwidth}{\tmplength}
}
\item[\textbf{Declaration}\hfill]
\ifpdf
\begin{flushleft}
\fi
\begin{ttfamily}
DefaultMorseSettings: TmorseSettings;\end{ttfamily}

\ifpdf
\end{flushleft}
\fi

\par
\item[\textbf{Description}]
A preset default of morse sounds at 18WPM, no noise, no Farnsworth Spacing

\end{list}
\chapter{Unit morse.dictionary.ascii}
\label{morse.dictionary.ascii}
\index{morse.dictionary.ascii}
\section{Description}
Contains the asciiToMorse array and setup for it from the MorseDictionary
\section{Variables}
\ifpdf
\subsection*{\large{\textbf{asciiToMorse}}\normalsize\hspace{1ex}\hrulefill}
\else
\subsection*{asciiToMorse}
\fi
\label{morse.dictionary.ascii-asciiToMorse}
\index{asciiToMorse}
\begin{list}{}{
\settowidth{\tmplength}{\textbf{Description}}
\setlength{\itemindent}{0cm}
\setlength{\listparindent}{0cm}
\setlength{\leftmargin}{\evensidemargin}
\addtolength{\leftmargin}{\tmplength}
\settowidth{\labelsep}{X}
\addtolength{\leftmargin}{\labelsep}
\setlength{\labelwidth}{\tmplength}
}
\item[\textbf{Declaration}\hfill]
\ifpdf
\begin{flushleft}
\fi
\begin{ttfamily}
asciiToMorse: array [0..127] of string;\end{ttfamily}

\ifpdf
\end{flushleft}
\fi

\par
\item[\textbf{Description}]
use asciiToMorse to quickly get the morse code string of an ascii character. Example \begin{ttfamily}asciiToMorse[Ord(a)]\end{ttfamily} shoud output a string of '.{-}', the dit{-}dah representing the letter A in morse code

\end{list}
\chapter{Unit morse.dictionary.util}
\label{morse.dictionary.util}
\index{morse.dictionary.util}
\section{Description}
A class Tdictionary{$<$}String, String{$>$} that returns the Value of the dit '.' or dah '{-}' of the given ASCII character Capital letters only! \texttt{\\\nopagebreak[3]
MorseLookup['A']\\
} Should return a string type with '.{-}'

The class is initialized and free when the unit is used in a program or library.
\section{Variables}
\ifpdf
\subsection*{\large{\textbf{MorseLookup}}\normalsize\hspace{1ex}\hrulefill}
\else
\subsection*{MorseLookup}
\fi
\label{morse.dictionary.util-MorseLookup}
\index{MorseLookup}
\begin{list}{}{
\settowidth{\tmplength}{\textbf{Description}}
\setlength{\itemindent}{0cm}
\setlength{\listparindent}{0cm}
\setlength{\leftmargin}{\evensidemargin}
\addtolength{\leftmargin}{\tmplength}
\settowidth{\labelsep}{X}
\addtolength{\leftmargin}{\labelsep}
\setlength{\labelwidth}{\tmplength}
}
\item[\textbf{Declaration}\hfill]
\ifpdf
\begin{flushleft}
\fi
\begin{ttfamily}
MorseLookup: TDictionary{$<$}string, string{$>$};\end{ttfamily}

\ifpdf
\end{flushleft}
\fi

\end{list}
\chapter{Unit morse.parser.Util}
\label{morse.parser.Util}
\index{morse.parser.Util}
\section{Overview}
\begin{description}
\item[\texttt{\begin{ttfamily}TmorseParser\end{ttfamily} Class}]
\end{description}
\section{Classes, Interfaces, Objects and Records}
\ifpdf
\subsection*{\large{\textbf{TmorseParser Class}}\normalsize\hspace{1ex}\hrulefill}
\else
\subsection*{TmorseParser Class}
\fi
\label{morse.parser.Util.TmorseParser}
\index{TmorseParser}
\subsubsection*{\large{\textbf{Hierarchy}}\normalsize\hspace{1ex}\hfill}
TmorseParser {$>$} TObject
%%%%Description
\subsubsection*{\large{\textbf{Methods}}\normalsize\hspace{1ex}\hfill}
\paragraph*{inputString}\hspace*{\fill}

\label{morse.parser.Util.TmorseParser-inputString}
\index{inputString}
\begin{list}{}{
\settowidth{\tmplength}{\textbf{Description}}
\setlength{\itemindent}{0cm}
\setlength{\listparindent}{0cm}
\setlength{\leftmargin}{\evensidemargin}
\addtolength{\leftmargin}{\tmplength}
\settowidth{\labelsep}{X}
\addtolength{\leftmargin}{\labelsep}
\setlength{\labelwidth}{\tmplength}
}
\item[\textbf{Declaration}\hfill]
\ifpdf
\begin{flushleft}
\fi
\begin{ttfamily}
public procedure inputString(const ATextInput: string);\end{ttfamily}

\ifpdf
\end{flushleft}
\fi

\end{list}
\paragraph*{getOutput}\hspace*{\fill}

\label{morse.parser.Util.TmorseParser-getOutput}
\index{getOutput}
\begin{list}{}{
\settowidth{\tmplength}{\textbf{Description}}
\setlength{\itemindent}{0cm}
\setlength{\listparindent}{0cm}
\setlength{\leftmargin}{\evensidemargin}
\addtolength{\leftmargin}{\tmplength}
\settowidth{\labelsep}{X}
\addtolength{\leftmargin}{\labelsep}
\setlength{\labelwidth}{\tmplength}
}
\item[\textbf{Declaration}\hfill]
\ifpdf
\begin{flushleft}
\fi
\begin{ttfamily}
public function getOutput: string;\end{ttfamily}

\ifpdf
\end{flushleft}
\fi

\end{list}
\chapter{Unit morse.pcm.pinknoise}
\label{morse.pcm.pinknoise}
\index{morse.pcm.pinknoise}
\section{Overview}
\begin{description}
\item[\texttt{PinkNoise}]
\end{description}
\section{Functions and Procedures}
\ifpdf
\subsection*{\large{\textbf{PinkNoise}}\normalsize\hspace{1ex}\hrulefill}
\else
\subsection*{PinkNoise}
\fi
\label{morse.pcm.pinknoise-PinkNoise}
\index{PinkNoise}
\begin{list}{}{
\settowidth{\tmplength}{\textbf{Description}}
\setlength{\itemindent}{0cm}
\setlength{\listparindent}{0cm}
\setlength{\leftmargin}{\evensidemargin}
\addtolength{\leftmargin}{\tmplength}
\settowidth{\labelsep}{X}
\addtolength{\leftmargin}{\labelsep}
\setlength{\labelwidth}{\tmplength}
}
\item[\textbf{Declaration}\hfill]
\ifpdf
\begin{flushleft}
\fi
\begin{ttfamily}
procedure PinkNoise(var aPCM: TNoise; aLengthms: UInt64; aAmplitude: UInt16 = 27000; aSampleRate: UInt32 = 44100);\end{ttfamily}

\ifpdf
\end{flushleft}
\fi

\par
\item[\textbf{Description}]
Generates pink noise (1/f noise) with specified parameters    \par
\item[\textbf{Parameters}]
\begin{description}
\item[aPCM] {-} Output array that will contain the generated samples
\item[aLengthms] {-} Duration of the noise in milliseconds
\item[aAmplitude] {-} Peak amplitude (default: 27000, max safe: 32767)
\item[aSampleRate] {-} Sample rate in Hz (default: 44100)
\end{description}


\end{list}
\section{Types}
\ifpdf
\subsection*{\large{\textbf{TNoise}}\normalsize\hspace{1ex}\hrulefill}
\else
\subsection*{TNoise}
\fi
\label{morse.pcm.pinknoise-TNoise}
\index{TNoise}
\begin{list}{}{
\settowidth{\tmplength}{\textbf{Description}}
\setlength{\itemindent}{0cm}
\setlength{\listparindent}{0cm}
\setlength{\leftmargin}{\evensidemargin}
\addtolength{\leftmargin}{\tmplength}
\settowidth{\labelsep}{X}
\addtolength{\leftmargin}{\labelsep}
\setlength{\labelwidth}{\tmplength}
}
\item[\textbf{Declaration}\hfill]
\ifpdf
\begin{flushleft}
\fi
\begin{ttfamily}
TNoise = array of Integer;\end{ttfamily}

\ifpdf
\end{flushleft}
\fi

\end{list}
\chapter{Unit morse.pcm.sine}
\label{morse.pcm.sine}
\index{morse.pcm.sine}
\section{Overview}
\begin{description}
\item[\texttt{sineWave}]
\item[\texttt{msToSamples}]
\end{description}
\section{Functions and Procedures}
\ifpdf
\subsection*{\large{\textbf{sineWave}}\normalsize\hspace{1ex}\hrulefill}
\else
\subsection*{sineWave}
\fi
\label{morse.pcm.sine-sineWave}
\index{sineWave}
\begin{list}{}{
\settowidth{\tmplength}{\textbf{Description}}
\setlength{\itemindent}{0cm}
\setlength{\listparindent}{0cm}
\setlength{\leftmargin}{\evensidemargin}
\addtolength{\leftmargin}{\tmplength}
\settowidth{\labelsep}{X}
\addtolength{\leftmargin}{\labelsep}
\setlength{\labelwidth}{\tmplength}
}
\item[\textbf{Declaration}\hfill]
\ifpdf
\begin{flushleft}
\fi
\begin{ttfamily}
procedure sineWave(var APcm : Tpcm; const length{\_}ms : Integer; const Hertz : Integer = 800; const Amplitude : Integer =1653562408; const sampleRate: Integer = 8000);\end{ttfamily}

\ifpdf
\end{flushleft}
\fi

\par
\item[\textbf{Description}]
type Tpcm = array of Integer;

\end{list}
\ifpdf
\subsection*{\large{\textbf{msToSamples}}\normalsize\hspace{1ex}\hrulefill}
\else
\subsection*{msToSamples}
\fi
\label{morse.pcm.sine-msToSamples}
\index{msToSamples}
\begin{list}{}{
\settowidth{\tmplength}{\textbf{Description}}
\setlength{\itemindent}{0cm}
\setlength{\listparindent}{0cm}
\setlength{\leftmargin}{\evensidemargin}
\addtolength{\leftmargin}{\tmplength}
\settowidth{\labelsep}{X}
\addtolength{\leftmargin}{\labelsep}
\setlength{\labelwidth}{\tmplength}
}
\item[\textbf{Declaration}\hfill]
\ifpdf
\begin{flushleft}
\fi
\begin{ttfamily}
function msToSamples(const ms: Integer; const sampleRate: Integer) : Integer; inline;\end{ttfamily}

\ifpdf
\end{flushleft}
\fi

\end{list}
\chapter{Unit morse.pcm.taper}
\label{morse.pcm.taper}
\index{morse.pcm.taper}
\section{Description}
This unit contains the code to taper the audio ends of the morse code pulse. Ends are tapered to 5 milli{-}seconds
\section{Overview}
\begin{description}
\item[\texttt{taper}]
\end{description}
\section{Functions and Procedures}
\ifpdf
\subsection*{\large{\textbf{taper}}\normalsize\hspace{1ex}\hrulefill}
\else
\subsection*{taper}
\fi
\label{morse.pcm.taper-taper}
\index{taper}
\begin{list}{}{
\settowidth{\tmplength}{\textbf{Description}}
\setlength{\itemindent}{0cm}
\setlength{\listparindent}{0cm}
\setlength{\leftmargin}{\evensidemargin}
\addtolength{\leftmargin}{\tmplength}
\settowidth{\labelsep}{X}
\addtolength{\leftmargin}{\labelsep}
\setlength{\labelwidth}{\tmplength}
}
\item[\textbf{Declaration}\hfill]
\ifpdf
\begin{flushleft}
\fi
\begin{ttfamily}
procedure taper(var Apcm: Tpcm; sampleRate : Integer = 8000);\end{ttfamily}

\ifpdf
\end{flushleft}
\fi

\par
\item[\textbf{Description}]
Pass the Tpcm and the sample rate. Default sample rate is 8000. Tpcm is modified. Do not send Tpcm with silence at end . Audio must start and stop at ends of Tpcm to be tapered properly.

\end{list}
\chapter{Unit morse.pcm.Util}
\label{morse.pcm.Util}
\index{morse.pcm.Util}
\section{Overview}
\begin{description}
\item[\texttt{\begin{ttfamily}TElementSettings\end{ttfamily} Record}]
\item[\texttt{\begin{ttfamily}TElementProperties\end{ttfamily} Class}]
\item[\texttt{\begin{ttfamily}TCreateElement\end{ttfamily} Class}]
\item[\texttt{\begin{ttfamily}TElementPCM\end{ttfamily} Class}]
\end{description}
\section{Classes, Interfaces, Objects and Records}
\ifpdf
\subsection*{\large{\textbf{TElementSettings Record}}\normalsize\hspace{1ex}\hrulefill}
\else
\subsection*{TElementSettings Record}
\fi
\label{morse.pcm.Util.TElementSettings}
\index{TElementSettings}
\subsubsection*{\large{\textbf{Description}}\normalsize\hspace{1ex}\hfill}
TNoiseSettings = record noiseType :TNoiseType; noiseLevel : single; //0.1 to 1 end;\subsubsection*{\large{\textbf{Fields}}\normalsize\hspace{1ex}\hfill}
\begin{list}{}{
\settowidth{\tmplength}{\textbf{MorseTimings}}
\setlength{\itemindent}{0cm}
\setlength{\listparindent}{0cm}
\setlength{\leftmargin}{\evensidemargin}
\addtolength{\leftmargin}{\tmplength}
\settowidth{\labelsep}{X}
\addtolength{\leftmargin}{\labelsep}
\setlength{\labelwidth}{\tmplength}
}
\label{morse.pcm.Util.TElementSettings-SampleRate}
\index{SampleRate}
\item[\textbf{SampleRate}\hfill]
\ifpdf
\begin{flushleft}
\fi
\begin{ttfamily}
public SampleRate: UInt32;\end{ttfamily}

\ifpdf
\end{flushleft}
\fi


\par  \label{morse.pcm.Util.TElementSettings-FreqHz}
\index{FreqHz}
\item[\textbf{FreqHz}\hfill]
\ifpdf
\begin{flushleft}
\fi
\begin{ttfamily}
public FreqHz: UInt32;\end{ttfamily}

\ifpdf
\end{flushleft}
\fi


\par  \label{morse.pcm.Util.TElementSettings-MorseTimings}
\index{MorseTimings}
\item[\textbf{MorseTimings}\hfill]
\ifpdf
\begin{flushleft}
\fi
\begin{ttfamily}
public MorseTimings: ImorseTiming;\end{ttfamily}

\ifpdf
\end{flushleft}
\fi


\par  \label{morse.pcm.Util.TElementSettings-noiseType}
\index{noiseType}
\item[\textbf{noiseType}\hfill]
\ifpdf
\begin{flushleft}
\fi
\begin{ttfamily}
public noiseType:TNoiseType;\end{ttfamily}

\ifpdf
\end{flushleft}
\fi


\par Noise : TnoiseSettings;\label{morse.pcm.Util.TElementSettings-noiseLevel}
\index{noiseLevel}
\item[\textbf{noiseLevel}\hfill]
\ifpdf
\begin{flushleft}
\fi
\begin{ttfamily}
public noiseLevel: single;\end{ttfamily}

\ifpdf
\end{flushleft}
\fi


\par  \end{list}
\ifpdf
\subsection*{\large{\textbf{TElementProperties Class}}\normalsize\hspace{1ex}\hrulefill}
\else
\subsection*{TElementProperties Class}
\fi
\label{morse.pcm.Util.TElementProperties}
\index{TElementProperties}
\subsubsection*{\large{\textbf{Hierarchy}}\normalsize\hspace{1ex}\hfill}
TElementProperties {$>$} TObject
\subsubsection*{\large{\textbf{Description}}\normalsize\hspace{1ex}\hfill}
maybe a class that calculates on creation?\subsubsection*{\large{\textbf{Properties}}\normalsize\hspace{1ex}\hfill}
\begin{list}{}{
\settowidth{\tmplength}{\textbf{sampleRate}}
\setlength{\itemindent}{0cm}
\setlength{\listparindent}{0cm}
\setlength{\leftmargin}{\evensidemargin}
\addtolength{\leftmargin}{\tmplength}
\settowidth{\labelsep}{X}
\addtolength{\leftmargin}{\labelsep}
\setlength{\labelwidth}{\tmplength}
}
\label{morse.pcm.Util.TElementProperties-sampleRate}
\index{sampleRate}
\item[\textbf{sampleRate}\hfill]
\ifpdf
\begin{flushleft}
\fi
\begin{ttfamily}
public property sampleRate : integer read fSampleRate write fSampleRate;\end{ttfamily}

\ifpdf
\end{flushleft}
\fi


\par fNoise : TnoiseSettings; fGenerator : TElementPCM\end{list}
\ifpdf
\subsection*{\large{\textbf{TCreateElement Class}}\normalsize\hspace{1ex}\hrulefill}
\else
\subsection*{TCreateElement Class}
\fi
\label{morse.pcm.Util.TCreateElement}
\index{TCreateElement}
\subsubsection*{\large{\textbf{Hierarchy}}\normalsize\hspace{1ex}\hfill}
TCreateElement {$>$} TObject
%%%%Description
\subsubsection*{\large{\textbf{Methods}}\normalsize\hspace{1ex}\hfill}
\paragraph*{create}\hspace*{\fill}

\label{morse.pcm.Util.TCreateElement-create}
\index{create}
\begin{list}{}{
\settowidth{\tmplength}{\textbf{Description}}
\setlength{\itemindent}{0cm}
\setlength{\listparindent}{0cm}
\setlength{\leftmargin}{\evensidemargin}
\addtolength{\leftmargin}{\tmplength}
\settowidth{\labelsep}{X}
\addtolength{\leftmargin}{\labelsep}
\setlength{\labelwidth}{\tmplength}
}
\item[\textbf{Declaration}\hfill]
\ifpdf
\begin{flushleft}
\fi
\begin{ttfamily}
public constructor create(settings : TElementSettings);\end{ttfamily}

\ifpdf
\end{flushleft}
\fi

\end{list}
\ifpdf
\subsection*{\large{\textbf{TElementPCM Class}}\normalsize\hspace{1ex}\hrulefill}
\else
\subsection*{TElementPCM Class}
\fi
\label{morse.pcm.Util.TElementPCM}
\index{TElementPCM}
\subsubsection*{\large{\textbf{Hierarchy}}\normalsize\hspace{1ex}\hfill}
TElementPCM {$>$} TObject
\subsubsection*{\large{\textbf{Description}}\normalsize\hspace{1ex}\hfill}
NOTE ========================== Submit Bug report for Blaise. It does not catch mismatch method declarations in class!

==============================\section{Types}
\ifpdf
\subsection*{\large{\textbf{EMorseElement}}\normalsize\hspace{1ex}\hrulefill}
\else
\subsection*{EMorseElement}
\fi
\label{morse.pcm.Util-EMorseElement}
\index{EMorseElement}
\begin{list}{}{
\settowidth{\tmplength}{\textbf{Description}}
\setlength{\itemindent}{0cm}
\setlength{\listparindent}{0cm}
\setlength{\leftmargin}{\evensidemargin}
\addtolength{\leftmargin}{\tmplength}
\settowidth{\labelsep}{X}
\addtolength{\leftmargin}{\labelsep}
\setlength{\labelwidth}{\tmplength}
}
\item[\textbf{Declaration}\hfill]
\ifpdf
\begin{flushleft}
\fi
\begin{ttfamily}
EMorseElement = Exception;\end{ttfamily}

\ifpdf
\end{flushleft}
\fi

\par
\item[\textbf{Description}]
, morse.pcm.pinknoise, legacey.random

\end{list}
\ifpdf
\subsection*{\large{\textbf{TNoiseType}}\normalsize\hspace{1ex}\hrulefill}
\else
\subsection*{TNoiseType}
\fi
\label{morse.pcm.Util-TNoiseType}
\index{TNoiseType}
\begin{list}{}{
\settowidth{\tmplength}{\textbf{Description}}
\setlength{\itemindent}{0cm}
\setlength{\listparindent}{0cm}
\setlength{\leftmargin}{\evensidemargin}
\addtolength{\leftmargin}{\tmplength}
\settowidth{\labelsep}{X}
\addtolength{\leftmargin}{\labelsep}
\setlength{\labelwidth}{\tmplength}
}
\item[\textbf{Declaration}\hfill]
\ifpdf
\begin{flushleft}
\fi
\begin{ttfamily}
TNoiseType = (...);\end{ttfamily}

\ifpdf
\end{flushleft}
\fi

\par
\item[\textbf{Description}]
TpcmArray{$<$}T{$>$} = array of T; //Not Accepted in Blaise yet!\item[\textbf{Values}]
\begin{description}
\item[\texttt{nsOff}] \label{morse.pcm.Util-nsOff}
\index{}
 
\item[\texttt{nsWhite}] \label{morse.pcm.Util-nsWhite}
\index{}
 
\item[\texttt{nsPink}] \label{morse.pcm.Util-nsPink}
\index{}
 
\item[\texttt{nsGaussien}] \label{morse.pcm.Util-nsGaussien}
\index{}
 
\end{description}


\end{list}
\chapter{Unit morse.pcm.whitenoise}
\label{morse.pcm.whitenoise}
\index{morse.pcm.whitenoise}
\section{Overview}
\begin{description}
\item[\texttt{whiteNoise}]
\end{description}
\section{Functions and Procedures}
\ifpdf
\subsection*{\large{\textbf{whiteNoise}}\normalsize\hspace{1ex}\hrulefill}
\else
\subsection*{whiteNoise}
\fi
\label{morse.pcm.whitenoise-whiteNoise}
\index{whiteNoise}
\begin{list}{}{
\settowidth{\tmplength}{\textbf{Description}}
\setlength{\itemindent}{0cm}
\setlength{\listparindent}{0cm}
\setlength{\leftmargin}{\evensidemargin}
\addtolength{\leftmargin}{\tmplength}
\settowidth{\labelsep}{X}
\addtolength{\leftmargin}{\labelsep}
\setlength{\labelwidth}{\tmplength}
}
\item[\textbf{Declaration}\hfill]
\ifpdf
\begin{flushleft}
\fi
\begin{ttfamily}
procedure whiteNoise(var aPCM: Tpcm; aLengthms: uint64; aAmplitude: Integer = 1653562408; aSampleRate: uint32 = 44100);\end{ttfamily}

\ifpdf
\end{flushleft}
\fi

\par
\item[\textbf{Description}]
type TNoise = array of int16;

\end{list}
\chapter{Program test}
\label{test}
\index{test}
\end{document}
